% GENERATED FILE. Edit canonical lessons, then run npm run generate:books.
% canonical-source-hash: fnv1a64:df1a8bae2c0bad84
% canonical-lessons: PT-C467-orgulho, PT-C467-vaidade, PT-C467-vaidoso, PT-C467-orgulhar, PT-C467-gabar

\chapter{Orgulho e vaidade}
\label{ch:pt-orgulho-e-vaidade}
\hlchaptermodality{467}

\begin{chapteropening}
\textbf{By the end of this chapter:} \emph{I can talk about pride, vanity and boasting.}

Complete the last lesson of chapter 467: I can talk about pride, vanity and boasting.
\end{chapteropening}

\section[o orgulho]{o orgulho — pride}

\label{lesson:PT-C467-orgulho}



\begin{quote}
Before the new one: say the Portuguese for a jungle, then the Portuguese for a slipper.
\end{quote}

\subsection*{You'll want to know: o orgulho}
\textbf{o orgulho} — \textquotedblleft{}pride\textquotedblright{}.

\begin{quote}
O meu pai fala do meu trabalho com muito orgulho.
\end{quote}

\begin{quote}
\emph{My father talks about my work with great pride.}
\end{quote}

\emph{Ter orgulho em alguém} (or \emph{de alguém}) — to be proud of someone.

\begin{cousinweb}
\emph{orgulho} came into Old Portuguese from Catalan \emph{orgull} \textquotedblleft{}pride\textquotedblright{} (EtymDB), which was borrowed from a Frankish word \emph{urgol} of the same meaning.
\end{cousinweb}

\begin{grammarlens}[title={What you've built}]
One more word for telling, explaining and describing at B1.
\end{grammarlens}

\subsection*{Your turn}
\begin{itemize}
\raggedright
  \item \emph{Say it:} \emph{o orgulho}
  \item \emph{Say it:} \emph{o orgulho}, once more
  \item \emph{Say it:} a sentence of your own with \emph{o orgulho}
  \item \emph{Recall:} say the Portuguese for a jungle, then the Portuguese for a slipper, then say \emph{o orgulho} again
\end{itemize}

\begin{tcolorbox}[breakable,colback=teal!4,colframe=teal!35!black,title={Before you move on}]
What does o orgulho mean? (\textquotedblleft{}Pride\textquotedblright{}.) Say it once more.
\end{tcolorbox}
\section[a vaidade]{a vaidade — vanity}

\label{lesson:PT-C467-vaidade}



\begin{quote}
Before the new one: say the Portuguese for a slipper, then the Portuguese for pride.
\end{quote}

\subsection*{You'll want to know: a vaidade}
\textbf{a vaidade} — \textquotedblleft{}vanity\textquotedblright{}.

\begin{quote}
Ela passa horas ao espelho por pura vaidade.
\end{quote}

\begin{quote}
\emph{She spends hours at the mirror out of pure vanity.}
\end{quote}

\emph{Por vaidade} — out of vanity.

\begin{cousinweb}
\emph{vaidade} is inherited, through Old Portuguese \emph{vãidade}, from Latin \emph{vānitās} \textquotedblleft{}emptiness\textquotedblright{} (EtymDB), built on \emph{vānus} \textquotedblleft{}empty\textquotedblright{}.
\end{cousinweb}

\begin{grammarlens}[title={What you've built}]
One more word for telling, explaining and describing at B1.
\end{grammarlens}

\subsection*{Your turn}
\begin{itemize}
\raggedright
  \item \emph{Say it:} \emph{a vaidade}
  \item \emph{Say it:} \emph{a vaidade}, once more
  \item \emph{Say it:} a sentence of your own with \emph{a vaidade}
  \item \emph{Recall:} say the Portuguese for a slipper, then the Portuguese for pride, then say \emph{a vaidade} again
\end{itemize}

\begin{tcolorbox}[breakable,colback=teal!4,colframe=teal!35!black,title={Before you move on}]
What does a vaidade mean? (\textquotedblleft{}Vanity\textquotedblright{}.) Say it once more.
\end{tcolorbox}
\section[vaidoso]{vaidoso — vain}

\label{lesson:PT-C467-vaidoso}



\begin{quote}
Before the new one: say the Portuguese for pride, then the Portuguese for vanity.
\end{quote}

\subsection*{You'll want to know: vaidoso}
\textbf{vaidoso} — \textquotedblleft{}vain\textquotedblright{}.

\begin{quote}
O meu irmão é vaidoso: compra sempre roupa cara.
\end{quote}

\begin{quote}
\emph{My brother is vain: he keeps buying expensive clothes.}
\end{quote}

Built on \emph{a vaidade}; a woman is \emph{vaidosa}.

\begin{cousinweb}
\emph{vaidoso} is built on \emph{vaidade} \textquotedblleft{}vanity\textquotedblright{} with \emph{-oso} \textquotedblleft{}full of\textquotedblright{}; \emph{vaidade} is inherited from Latin \emph{vānitās} \textquotedblleft{}emptiness\textquotedblright{} (EtymDB).
\end{cousinweb}

\begin{grammarlens}[title={What you've built}]
One more word for telling, explaining and describing at B1.
\end{grammarlens}

\subsection*{Your turn}
\begin{itemize}
\raggedright
  \item \emph{Say it:} \emph{vaidoso}
  \item \emph{Say it:} \emph{vaidoso}, once more
  \item \emph{Say it:} a sentence of your own with \emph{vaidoso}
  \item \emph{Recall:} say the Portuguese for pride, then the Portuguese for vanity, then say \emph{vaidoso} again
\end{itemize}

\begin{tcolorbox}[breakable,colback=teal!4,colframe=teal!35!black,title={Before you move on}]
What does vaidoso mean? (\textquotedblleft{}Vain\textquotedblright{}.) Say it once more.
\end{tcolorbox}
\section[orgulhar]{orgulhar — to be proud (orgulhar-se de); to make proud}

\label{lesson:PT-C467-orgulhar}



\begin{quote}
Before the new one: say the Portuguese for vanity, then the Portuguese for vain.
\end{quote}

\subsection*{You'll want to know: orgulhar}
\textbf{orgulhar} — \textquotedblleft{}to be proud (orgulhar-se de); to make proud\textquotedblright{}.

\begin{quote}
A avó orgulha-se muito dos netos.
\end{quote}

\begin{quote}
\emph{Grandma is very proud of her grandchildren.}
\end{quote}

Mostly reflexive, with the pronoun after the verb: \emph{orgulho-me de ti} — I am proud of you. In Brazil: \emph{me orgulho}.

\begin{cousinweb}
\emph{orgulhar} is built on \emph{orgulho} \textquotedblleft{}pride\textquotedblright{}, which came into Portuguese from Catalan \emph{orgull} (EtymDB), a borrowing from Frankish.
\end{cousinweb}

\begin{grammarlens}[title={What you've built}]
One more word for telling, explaining and describing at B1.
\end{grammarlens}

\subsection*{Your turn}
\begin{itemize}
\raggedright
  \item \emph{Say it:} \emph{orgulhar}
  \item \emph{Say it:} \emph{orgulhar}, once more
  \item \emph{Say it:} a sentence of your own with \emph{orgulhar}
  \item \emph{Recall:} say the Portuguese for vanity, then the Portuguese for vain, then say \emph{orgulhar} again
\end{itemize}

\begin{tcolorbox}[breakable,colback=teal!4,colframe=teal!35!black,title={Before you move on}]
What does orgulhar mean? (\textquotedblleft{}To be proud (orgulhar-se de); to make proud\textquotedblright{}.) Say it once more.
\end{tcolorbox}
\section[gabar]{gabar — to boast (gabar-se); to praise}

\label{lesson:PT-C467-gabar}



\begin{quote}
Before the new one: say the Portuguese for vain, then the Portuguese for to be proud.
\end{quote}

\subsection*{You'll want to know: gabar}
\textbf{gabar} — \textquotedblleft{}to boast (gabar-se); to praise\textquotedblright{}.

\begin{quote}
O meu colega gaba-se sempre do carro novo.
\end{quote}

\begin{quote}
\emph{My colleague is forever boasting about his new car.}
\end{quote}

Mostly reflexive, with the pronoun after the verb: \emph{gabar-se de} — to boast about. Without \emph{-se}, \emph{gabar alguém} is to praise someone, a less common use.

\begin{cousinweb}
\emph{gabar} comes from Occitan \emph{gabar} \textquotedblleft{}to boast\textquotedblright{} (EtymDB), which EtymDB traces to Old Norse \emph{gabb}, a word usually glossed \textquotedblleft{}mockery\textquotedblright{}.
\end{cousinweb}

\begin{grammarlens}[title={What you've built}]
One more word for telling, explaining and describing at B1.
\end{grammarlens}

\subsection*{Your turn}
\begin{itemize}
\raggedright
  \item \emph{Say it:} \emph{gabar}
  \item \emph{Say it:} \emph{gabar}, once more
  \item \emph{Say it:} a sentence of your own with \emph{gabar}
  \item \emph{Recall:} say the Portuguese for vain, then the Portuguese for to be proud, then say \emph{gabar} again
\end{itemize}

\begin{tcolorbox}[breakable,colback=teal!4,colframe=teal!35!black,title={Before you move on}]
What does gabar mean? (\textquotedblleft{}To boast (gabar-se); to praise\textquotedblright{}.) Say it once more.
\end{tcolorbox}
